\section{Memory Compression and Startup Covariance}
\newcommand{\MemMean}{25}
\newcommand{\MemMoments}{43.5}
\newcommand{\MemMixture}{45.1}
\newcommand{\CovIIDScale}{0.0049}
\newcommand{\CovIIDSample}{2.52}
\newcommand{\CovIIDShrink}{0.00671}
\newcommand{\CovARScale}{0.0391}
\newcommand{\CovARSample}{42}
\newcommand{\CovARShrink}{20.4}


\subsection{Four Matched Representations of Recovery}
An external controlled example uses three swing nodes, a grounded
reference, two possible recovery components and three measured coordinates.
It compares a state mean, a moment-matched Gaussian, the full state
mixture and the same mixture with historical cause labels. All arms use
the same linear dynamics and class priors. The next load can occur at
either of two sources at a known onset. Each arm uses 3,000 independent
normal windows to set a 5\% window false-positive target. Each scenario
then evaluates 3,800 new windows. This target concerns fixed windows,
whereas the article calibrates alarms per hour of continuous monitoring.

In the nominal colored-noise scenario, the mean, Gaussian and mixture
correctly detect and locate \MemMean{}\%, \MemMoments{}\% and
\MemMixture{}\% of event windows, respectively. Across the four
scenarios, the mixture and labeled histories have identical new-source
posteriors. This equality follows from the shared likelihoods and
transition law after summing over component labels. Labels support
queries about the previous cause. The recursive Gaussian calculation
agrees with the batch calculation on every prefix of the checked
trajectory. These results support the distinction in the predictive
sufficiency analysis. The numerical gains refer to the small network and
its prescribed recovery distribution.

\subsection{Accuracy of a Startup Covariance}
For true and estimated positive definite covariances $\Sigma$ and
$\widehat\Sigma$, let $\eta=\|\Sigma^{-1/2}\widehat\Sigma
\Sigma^{-1/2}-I\|_2$. When $\eta<1$, every fixed signature $d$ satisfies
\begin{equation}
\frac{d^\top\Sigma^{-1}d}{1+\eta}
\le d^\top\widehat\Sigma^{-1}d
\le\frac{d^\top\Sigma^{-1}d}{1-\eta}.
\end{equation}
This follows by inverting the eigenvalue bounds in whitened coordinates.
The inequality controls a fixed quadratic. A data-driven alarm rule also
depends on log determinants, selection and the normal predictive law.
The observed TVE is on a signal scale, whereas $\eta$ measures relative
covariance error.

We generated 128 Gaussian prefixes per setting in 47 coordinates, using
the nominal innovation correlation as the reference. The shifted
condition adds unit diagonal variance in standardized coordinates.
At 300 reports in that condition, median excess marginal log loss in
nats per vector is \CovIIDScale{}, \CovIIDSample{} and \CovIIDShrink{}
for a scaled nominal covariance, a sample covariance and Ledoit--Wolf
shrinkage, respectively, with independent draws. For an AR(1) coefficient
of 0.9 the corresponding losses are \CovARScale{}, \CovARSample{} and
\CovARShrink{}. Shrinkage improves on the unstructured sample estimate
in these settings. The nominal correlation remains a useful prior in
this construction. Recovering load attribution at physical independent
PMU TVE requires testing the covariance in the full bank with independent
calibration and evaluation.
